\section{Summary of Findings}
This thesis examined two linked questions in EEG-based dementia classification. First, we asked whether a correct first-minute prediction remains consistent during a longer recording. Second, we asked whether a custom decoder attached to a pretrained EEG encoder is competitive when participants, rather than recording windows, are used for evaluation. The two investigations share a concern about reliability but use different analytic cohorts and experimental protocols.

In the original temporal experiment, an RBF SVM achieved first-minute balanced accuracy of 73.48\% for AD versus controls and 73.72\% for FTD versus controls. Reported predictions changed during the five-minute evaluation for 44.35\% and 33.66\% of cases, respectively. Stable but wrong decisions were also present. A separate MCI-versus-control study showed that low label-change rates can coexist with poor discrimination. These findings support the principle that prediction stability and prediction correctness must be evaluated together.

The later EEGraft architecture used a frozen CBraMod EEG encoder with a custom dual-scale residual fusion decoder. On 52 development participants, it achieved 77.06\% pooled balanced accuracy. XGBoost was very close at 76.61\%, and both models correctly classified 40 of the 52 participants. On the locked group of 13 participants, DSRF achieved 70.24\% balanced accuracy and XGBoost achieved 77.38\%. The paired uncertainty interval crossed zero, so superiority of either architecture in a wider population was not statistically established. The one-time observed holdout result nevertheless favored XGBoost.

Artificial noise and single-channel blanking affected DSRF accuracy, label consistency and probability scores. Electrode-masking results exposed model input sensitivity but did not establish neural causes or biomarkers. The overall result is a carefully bounded internal study rather than a clinically validated diagnostic system.

\section{Contributions}
The main contributions are: (1) a subject-level temporal analysis that separates changing predictions from stable errors; (2) an implemented gated dual-scale residual fusion decoder using frozen pretrained representations and a learned temporal pathway; (3) controlled development comparisons with several classical and neural alternatives, including an ablation of the residual pathway; (4) a one-time held-out assessment that retains the unfavorable result for the proposed decoder; and (5) robustness, calibration and input-masking analyses that show why one score is not enough. The experiments do not demonstrate an original foundation model or a confirmed new clinical biomarker.

\section{Recommendations for Future Work}
A matched experiment should apply the same five-minute temporal protocol to DSRF, XGBoost and the original spectral SVM on a reconciled subject set. This would directly test whether the custom decoder improves temporal consistency without sacrificing accuracy. The study should then extend to independent AD recordings acquired in different institutions, using a predetermined channel-mapping and preprocessing rule. This would address the main gap in external validity.

Larger independently labeled cohorts are needed for reliable estimates of patient-level performance and probability calibration. Multiple training seeds and uncertainty estimates that include retraining variation would improve architecture comparisons. More realistic artifact and missing-channel experiments should use physically plausible disturbances, not only synthetic Gaussian noise or a zeroed channel. Finally, clinicians and EEG experts should evaluate whether model explanations remain meaningful under genuine acquisition differences. These steps are proposed future work; none are claimed to have been completed here.

\section{Final Conclusion}
EEG-based classification can provide useful research signals while still producing inconsistent or persistently wrong individual predictions. EEGraft demonstrates a viable custom-decoder implementation and a transparent way to examine its performance alongside traditional models. Its reported results support continued methodological investigation, not clinical diagnostic use or future-risk prediction. The central lesson is that reliability must be evaluated as carefully as accuracy, and that negative or inconclusive results are necessary parts of a credible scientific conclusion.
